\documentclass[12pt]{scrartcl}
\usepackage{helvethia} % eksperimen tema Helvethia; kembali ke tampilan lama: \usepackage[hagavi]{azzam}
\providecommand{\sumber}[1]{(#1)}

% Fungsi Floor, Ceiling, dan Fractional Part - Modul (gabungan Materi dan Soal Latihan)
% Materi diadaptasi dari `1 Aljabar/Lantai dan Atap.tex`:
% definisi floor, ceiling, dan fractional part beserta sifat dasarnya
% (ditulis ulang dengan contoh pengerjaan tiap blok teori).
% Soal Latihan dari `1 Aljabar/Lantai dan Atap.tex`: soal OSK 2013, 2016, 2018,
% 2019, 2021, JBMO 2021 (dimodifikasi), Hong Kong Prelim 1999, Putnam 1986,
% USAMO 1981, IMO 1968, dan soal tanpa sumber.
% _Solusi tetap terpisah.
% Sumber tiap soal latihan dimunculkan di PDF sebagai \item (Sumber); komentar menyimpan detail lengkap.

\title{Fungsi Floor, Ceiling, dan Fractional Part}
\subtitle{Modul Aljabar}
\date{20 November 2026}

\begin{document}
\maketitle

\renewcommand*\contentsname{Daftar Isi}
\tableofcontents

\newpage

\part*{Materi}
\section{Floor, Ceiling, dan Fractional Part}

\begin{definition}
Untuk $x \in \RR$, \textbf{floor} dari $x$, ditulis $\floor{x}$,
adalah bilangan bulat terbesar yang kurang dari atau sama dengan $x$.
\textbf{Ceiling} dari $x$, ditulis $\ceiling{x}$,
adalah bilangan bulat terkecil yang lebih dari atau sama dengan $x$.
\istilah{Bagian pecahan}{fractional part} dari $x$, ditulis $\{x\}$,
didefinisikan sebagai
\[ \{x\} = x - \floor{x}. \]
\end{definition}

\begin{example}
Hitung $\floor{\pi}$, $\ceiling{5}$, dan $\{-7/2\}$.
\end{example}
\begin{proof}[Solusi]
Bilangan bulat terbesar yang tidak melebihi $\pi \approx 3{,}14$ adalah $3$:
\[ \floor{\pi} = 3. \]
Karena $5$ sudah bulat, ceiling-nya dirinya sendiri:
\[ \ceiling{5} = 5. \]
Untuk $-7/2 = -3{,}5$, bilangan bulat terbesar di bawahnya adalah $-4$:
\[ \floor{-7/2} = -4. \]
Akibatnya
\[ \{-7/2\} = -\frac{7}{2} - (-4) = \frac{1}{2}. \]
\end{proof}

\begin{fact}
Untuk setiap $x \in \RR$ berlaku
\[ \floor{x} \le x < \floor{x} + 1
\quad\text{dan}\quad
\ceiling{x} - 1 < x \le \ceiling{x}, \]
serta $0 \le \{x\} < 1$.
\end{fact}

\begin{example}
Jika $\floor{x} = 5$, tentukan semua nilai $x$ yang mungkin.
\end{example}
\begin{proof}[Solusi]
Dari ketaksamaan dasar floor,
\[ 5 \le x < 5 + 1. \]
Jadi $x \in [5, 6)$.
\end{proof}

\section{Sifat-Sifat Dasar}

\begin{fact}
Untuk $x, y \in \RR$ dan $n \in \ZZ$ berlaku:
\begin{enumerate}
    \item $\floor{n + x} = n + \floor{x}$,
    $\ceiling{n + x} = n + \ceiling{x}$, dan $\{n + x\} = \{x\}$.
    \item Jika $x \le y$ maka $\floor{x} \le \floor{y}$
    dan $\ceiling{x} \le \ceiling{y}$.
    \item $\floor{x} = \ceiling{x}$ jika dan hanya jika $x \in \ZZ$;
    jika $x \notin \ZZ$ maka $\ceiling{x} = \floor{x} + 1$.
    \item $\floor{x} + \floor{y} \le \floor{x + y}
    \le \floor{x} + \floor{y} + 1$.
\end{enumerate}
\end{fact}

\begin{example}
Hitung $\floor{100 + \sqrt{2}}$.
\end{example}
\begin{proof}[Solusi]
Karena $1 < \sqrt{2} < 2$, didapat $\floor{\sqrt{2}} = 1$.
Geser keluar bilangan bulatnya:
\[ \floor{100 + \sqrt{2}} = 100 + \floor{\sqrt{2}} = 101. \]
\end{proof}

\begin{example}
Buktikan bahwa untuk setiap $x \in \RR$ berlaku
\[ \floor{x} + \floor{-x} =
\begin{cases}
0, & x \in \ZZ, \\
-1, & x \notin \ZZ.
\end{cases} \]
\end{example}
\begin{proof}[Solusi]
Jika $x = n$ bulat, maka $\floor{x} + \floor{-x} = n + (-n) = 0$.
Jika $x$ tidak bulat, tulis $x = n + f$ dengan $n \in \ZZ$
dan $0 < f < 1$. Maka
\begin{align*}
\floor{x} + \floor{-x} &= n + \floor{-n - f} \\
&= n + (-n - 1) \\
&= -1.
\end{align*}
\end{proof}

\section{Persamaan Floor}

\begin{theorem}[Karakterisasi floor]
Untuk $x \in \RR$ dan $n \in \ZZ$ berlaku
\[ \floor{x} = n \quad\text{jika dan hanya jika}\quad n \le x < n + 1. \]
\end{theorem}

Substitusi $n = \floor{x}$ mengubah persamaan floor menjadi
ketaksamaan biasa dalam $x$. Selesaikan ketaksamaannya, lalu periksa kembali.

\begin{example}
Selesaikan $\floor{3x + 1} = 7$ untuk $x \in \RR$.
\end{example}
\begin{proof}[Solusi]
Menurut karakterisasi floor,
\[ 7 \le 3x + 1 < 8. \]
Kurangkan $1$ lalu bagi dengan $3$:
\[ 2 \le x < \frac{7}{3}. \]
\end{proof}

\begin{example}
Selesaikan $\floor{x}^2 - 5\floor{x} + 6 = 0$ untuk $x \in \RR$.
\end{example}
\begin{proof}[Solusi]
Tulis $n = \floor{x}$. Persamaannya menjadi
\[ n^2 - 5n + 6 = (n - 2)(n - 3) = 0, \]
jadi $n = 2$ atau $n = 3$.
Untuk $n = 2$ diperoleh $2 \le x < 3$,
untuk $n = 3$ diperoleh $3 \le x < 4$.
Gabungannya $x \in [2, 4)$.
\end{proof}

\section{Identitas Klasik}

\begin{lemma}
Untuk setiap $x \in \RR$ berlaku
\[ \floor{x + \tfrac12} = \floor{2x} - \floor{x}. \]
\end{lemma}
\begin{proof}
Tulis $x = n + f$ dengan $n \in \ZZ$ dan $0 \le f < 1$.
Jika $f < \frac12$, maka
\[ \floor{x + \tfrac12} = n
\quad\text{dan}\quad
\floor{2x} - \floor{x} = 2n - n = n. \]
Jika $f \ge \frac12$, maka
\[ \floor{x + \tfrac12} = n + 1
\quad\text{dan}\quad
\floor{2x} - \floor{x} = (2n + 1) - n = n + 1. \]
\end{proof}

\begin{example}
Periksa kesamaan lemma di atas untuk $x = 11/4$.
\end{example}
\begin{proof}[Solusi]
Ruas kirinya
\[ \floor{11/4 + 1/2} = \floor{13/4} = 3. \]
Ruas kanannya
\[ \floor{11/2} - \floor{11/4} = 5 - 2 = 3. \]
Keduanya sama.
\end{proof}

\begin{theorem}[Identitas Hermite]
Untuk $x \in \RR$ dan bilangan asli $n$ berlaku
\[ \floor{nx} = \floor{x} + \floor{x + \tfrac1n}
+ \floor{x + \tfrac2n} + \dots + \floor{x + \tfrac{n-1}{n}}. \]
\end{theorem}
\begin{proof}
Cukup membuktikan untuk $0 \le x < 1$,
karena kedua ruas bergeser sama jauhnya oleh bilangan bulat.
Tulis $m = \floor{nx}$, jadi
\[ \frac{m}{n} \le x < \frac{m+1}{n}. \]
Untuk $k = 0, 1, \dots, n-1$, suku $\floor{x + k/n}$
bernilai $0$ saat $x + k/n < 1$ dan $1$ saat $1 \le x + k/n < 2$.
Batasnya di $k = n - m$: untuk $k \le n - m - 1$ berlaku
\[ x + \frac{k}{n} < \frac{m+1}{n} + \frac{n-m-1}{n} = 1, \]
sedangkan untuk $k \ge n - m$ berlaku
\[ x + \frac{k}{n} \ge \frac{m}{n} + \frac{n-m}{n} = 1. \]
Tepat $m$ suku bernilai $1$. Jumlahnya $m = \floor{nx}$.
\end{proof}

\begin{example}
Hitung $\floor{12/5} + \floor{12/5 + 1/3} + \floor{12/5 + 2/3}$.
\end{example}
\begin{proof}[Solusi]
Menurut Identitas Hermite dengan $x = 12/5$ dan $n = 3$,
jumlahnya sama dengan
\[ \floor{3 \cdot 12/5} = \floor{36/5} = 7. \]
Pemeriksaan langsung: $2 + 2 + 3 = 7$.
\end{proof}

\section{Soal Latihan}

\subsection{Dasar-Dasar Floor dan Ceiling}
\begin{enumerate}

	% Source: Lantai dan Atap.tex (Latihan Soal 1, diterjemahkan ke Bahasa Indonesia, notasi diseragamkan)
	% Jawaban: 4, yaitu n = 6, 12, 18, 24
	\item
	Misalkan $\floor{x}$ menyatakan bilangan bulat terbesar yang tidak melebihi $x$.
	Ada berapa banyak bilangan asli $n$ sehingga
	\[ \floor{\frac{n}{5}} = \frac{n}{6} ? \]

	% Source: Lantai dan Atap.tex (Latihan Soal 2, diterjemahkan ke Bahasa Indonesia)
	% Jawaban: x = 10/3
	\item
	Jika $\floor{x}\{x\} = 1$ dan
	$\floor{x}^2 - \floor{x}(1+x) + 4 = 0$,
	tentukan nilai $x$.

	% Source: OSK 2013 (esai, dari Lantai dan Atap.tex Latihan Soal 4)
	% Jawaban: 2 < x < 3
	\item \sumber{OSK 2013}
	Tentukan semua $x$ yang memenuhi
	\[ \floor{x} + \ceiling{x} = 5. \]

\end{enumerate}

\subsection{Persamaan Floor}
\begin{enumerate}[resume]

	% Source: Lantai dan Atap.tex (Latihan Soal 3, diterjemahkan ke Bahasa Indonesia)
	% Jawaban: 55/7
	\item
	Tentukan bilangan real positif terkecil $x$ sehingga
	\[ \floor{x^2} - x\floor{x} = 6. \]

	% Source: OSK 2016 (dari Lantai dan Atap.tex Latihan Soal 5)
	% Jawaban: 516
	\item \sumber{OSK 2016}
	Banyaknya bilangan asli $n \in \{1, 2, 3, \dots, 1000\}$
	sehingga terdapat bilangan real positif $x$ yang memenuhi
	$x^2 + \floor{x}^2 = n$ adalah \dots

	% Source: OSK 2018 (dari Lantai dan Atap.tex Latihan Soal 6)
	% Jawaban: 274
	\item \sumber{OSK 2018}
	Jika diketahui $\floor{x} + \floor{y} + y = 43{,}8$ dan
	$x + y - \floor{x} = 18{,}4$,
	nilai $10(x + y)$ adalah \dots

	% Source: OSK 2019 (dari Lantai dan Atap.tex Latihan Soal 7)
	% Jawaban: 3
	\item \sumber{OSK 2019}
	Interval $[a, b)$ adalah himpunan semua bilangan real $x$ yang memenuhi
	\[ \floor{2x}^2 = \ceiling{x} + 7. \]
	Nilai $a \cdot b$ adalah \dots

	% Source: osk sma 2021 nomor 1 isian singkat (dari Lantai dan Atap.tex Latihan Soal 8;
	% berkas sumber salah tulis jumlah solusinya sebagai p generik, dikembalikan ke 51
	% sesuai naskah OSK 2021 asli)
	% Jawaban: a = 25
	\item \sumber{OSK 2021}
	Jika $a > 1$ suatu bilangan asli sehingga hasil penjumlahan semua
	bilangan real $x$ yang memenuhi persamaan
	\[ \floor{x}^2 - 2ax + a = 0 \]
	adalah $51$, maka $a$ adalah \dots

	% Source: Lantai dan Atap.tex (Latihan Soal 9, redaksi dirapikan)
	% Jawaban: 2, yaitu x = -3/4 dan x = 19/4
	\item
	Carilah banyaknya bilangan real $x$ yang memenuhi
	\[ \floor{x^2} = 4x + 3. \]

	% Source: Lantai dan Atap.tex (Latihan Soal 10, redaksi dirapikan)
	% Jawaban: 1215
	\item
	Jika $x$ adalah suatu bilangan real yang memenuhi persamaan
	\[ \floor{x} + \floor{2x} + \floor{3x} + \floor{4x} = 2024, \]
	tentukan nilai dari $\floor{6x}$.

	% Source: Lantai dan Atap.tex (Latihan Soal 11, redaksi dirapikan)
	% Jawaban: 232
	\item
	Jumlah dari semua solusi bilangan real $x$ dari persamaan
	\[ 2\floor{x}^2 + 3\{x\}^2 = \frac{7}{4}x\floor{x} \]
	dapat dinyatakan dalam bentuk $\frac{p}{q}$ dengan $p$ dan $q$
	bilangan bulat yang saling relatif prima.
	Tentukan nilai dari $10p + q$.

	% Source: Lantai dan Atap.tex (Latihan Soal 12, modifikasi JBMO 2021)
	% Jawaban: x = 5/36 dan x = 1/9
	\item \sumber{JBMO 2021}
	Carilah seluruh penyelesaian dari persamaan
	\[ 2\floor{\frac{1}{2x}} - 7 = 9(1 - 8x). \]

\end{enumerate}

\subsection{Identitas dan Jumlah Floor}
\begin{enumerate}[resume]

	% Source: Hong Kong IMO Prelim 1999-2000 (dari Lantai dan Atap.tex Sepuh 3,
	% diterjemahkan ke Bahasa Indonesia)
	% Jawaban: 1165
	\item \sumber{Hong Kong Prelim 1999}
	Tentukan bilangan bulat $n$ yang memenuhi
	\[ \floor{\frac{n}{1!}} + \floor{\frac{n}{2!}}
	+ \dots + \floor{\frac{n}{10!}} = 1999. \]

	% Source: Putnam 1986 (dari Lantai dan Atap.tex Sepuh 4, diterjemahkan ke Bahasa Indonesia)
	% Jawaban: 3
	\item \sumber{Putnam 1986}
	Berapakah digit satuan dari
	\[ \floor{\frac{10^{20000}}{10^{100} + 3}} ? \]

	% Source: USAMO 1981 (dari Lantai dan Atap.tex Sepuh 5, diterjemahkan ke Bahasa Indonesia)
	\item \sumber{USAMO 1981}
	Jika $x$ bilangan real positif dan $n$ bilangan asli, buktikan bahwa
	\[ \floor{nx} \ge \frac{\floor{x}}{1} + \frac{\floor{2x}}{2}
	+ \frac{\floor{3x}}{3} + \dots + \frac{\floor{nx}}{n}. \]

	% Source: IMO 1968 (dari Lantai dan Atap.tex Sepuh 6, diterjemahkan ke Bahasa Indonesia)
	% Jawaban: n
	\item \sumber{IMO 1968}
	Jika $n$ bilangan asli, nyatakan sebagai fungsi sederhana dari $n$ jumlah
	\[ \floor{\frac{n+1}{2}} + \floor{\frac{n+2}{4}}
	+ \dots + \floor{\frac{n+2^k}{2^{k+1}}} + \dots \]

	% Source: Lantai dan Atap.tex (Sepuh 7, tanpa sumber luar; redaksi dirapikan)
	\item
	Misalkan $n$ bilangan asli. Buktikan bahwa
	\[ \floor{\frac{n}{1}} + \floor{\frac{n}{2}} + \dots
	+ \floor{\frac{n}{n}} + \floor{\sqrt{n}} \]
	merupakan bilangan genap.

	% Source: Lantai dan Atap.tex (Sepuh 8, tanpa sumber luar; diterjemahkan ke Bahasa Indonesia)
	% Jawaban: semua n dengan n + 1 bilangan kuadrat
	\item
	Definisikan $q(n) = \floor{n/\floor{\sqrt{n}}}$ untuk $n \in \NN$.
	Tentukan semua $n$ sehingga $q(n) > q(n+1)$.

	% Source: Lantai dan Atap.tex (Sepuh 9, tanpa sumber luar; diterjemahkan ke Bahasa Indonesia)
	\item
	Misalkan $r \ge 1$ bilangan real sehingga untuk setiap bilangan asli
	$m$ dan $n$ dengan $m \mid n$ berlaku pula
	$\floor{mr} \mid \floor{nr}$.
	Buktikan bahwa $r$ bilangan bulat.

\end{enumerate}

\end{document}
